% !TITLE 春日忆李白
% !AUTHOR 杜甫
% !DYNASTY 唐
% !ORDER 32

\begin{poem}
白也诗无敌\textsuperscript{1}，飘然思不群\textsuperscript{2}。
清新庾开府\textsuperscript{3}，俊逸鲍参军\textsuperscript{4}。
渭北春天树，江东日暮云\textsuperscript{5}。
何时一樽酒，重与细论文\textsuperscript{6}。
\end{poem}

\begin{annotations}
\notetext{1}{白也诗无敌：李白的诗举世无敌。“白也”——直接称呼李白的名，像是朋友间的赞叹。}
\notetext{2}{飘然思不群：才思飘逸高远，超出流俗，与众不同。}
\notetext{3}{庾开府：庾信，南北朝时期著名文学家，诗风清新。“开府”是他的官名（开府仪同三司）。}
\notetext{4}{鲍参军：鲍照，南朝宋诗人，诗风俊逸。“参军”是他的官名（前军参军）。杜甫用这两位前代诗人来比拟李白。}
\notetext{5}{渭北春天树，江东日暮云：我在渭北（杜甫在长安）看着春天的树，你在江东（李白在江南）看着日暮的云。相隔千里，各看各的风景，心里却想着对方。}
\notetext{6}{细论文：细细地讨论诗文。什么时候能再喝一杯酒，重温论诗的旧梦？这个愿望杜甫终生未能实现。}
\end{annotations}

\begin{appreciation}
这是杜甫写给李白的诗。天宝五载（746年），杜甫与李白在鲁郡分别，此后两人再也没有见过面。但杜甫一生都在思念李白，为他写了近二十首诗，这是其中最动人的一首。

白也诗无敌，飘然思不群——李白的诗举世无敌，他的才思飘逸，与众不同。白也这个称呼很亲切，像是一个朋友在背后赞叹另一个朋友。无敌不是夸张，是杜甫发自内心的判断。飘然思不群——飘然是形容李白的诗风，也是形容他的为人：不受拘束，天马行空。

清新庾开府，俊逸鲍参军——他的诗清新得像庾信，俊逸得像鲍照。庾开府就是庾信，鲍参军就是鲍照，都是南北朝时期的著名诗人。杜甫用这两个古人来比拟李白，是在说：李白的才华不仅超越了当世，甚至可以和历史上最伟大的诗人比肩。

渭北春天树，江东日暮云——我在渭北，看着春天的树；你在江东，看着日暮的云。这是全诗最动人的两句。渭北是杜甫所在的地方，江东是李白所在的地方。两人相隔千里，各看各的风景。但杜甫说：我看树的时候，一定也想着你看云的样子。这种遥隔千里的相互凝视，是朋友之间最深沉的牵挂。春天树和日暮云，一春一暮，一树一云，对仗工整，却充满了柔情。

何时一樽酒，重与细论文——什么时候能再一起喝一杯酒，细细地讨论诗歌呢？这个愿望，杜甫终生未能实现。李白不久后死于当涂，杜甫直到去世，都没有再见到他。但这首诗让我们看到：真正的友谊，不需要朝夕相处，只需要彼此懂得。杜甫懂李白，懂到可以用二十个字概括他的全部才华；这种懂得，比一千次见面都更珍贵。
\end{appreciation}
